% ---------------------------------------------------------------------------
% The two specific cases K(2k+1, k) and K(2k+2, k).
% Transcribed from handwritten notes (pages 3 to 6).
% Every departure from the handwritten wording is marked with a comment
% beginning "% CHANGED:" on the line above it.
% ---------------------------------------------------------------------------
\documentclass[11pt]{article}
\usepackage[margin=1in]{geometry}
\usepackage{amsmath,amssymb}
\setlength{\parindent}{0pt}
\setlength{\parskip}{0.9em}

\title{Kneser graphs $K(2k+n,\,k)$ with $n>0$ are not CFON$^{*}$ 1-colorable:\\[2pt]
       the cases $K(2k+1,\,k)$ and $K(2k+2,\,k)$}
\author{Shrey Shah}
\date{}

\begin{document}
\maketitle

% CHANGED: this opening is added so the file can be read on its own. The
% wording is adapted from pages 1, 2 and 5 of the notes. The definition of an
% efficient open dominating set is not repeated here.
Let $K$ be the Kneser graph being considered, and let $T$ be an efficient open
dominating set of $K$, the set of all colored vertices in a CFON$^{*}$
1-coloring.

Let each vertex, $v$, in $K$ be assigned a variable:
\[
  x_v =
  \begin{cases}
    1, & v \in T\\
    0, & v \notin T
  \end{cases}
\]

Let $A$ be the adjacency matrix, with each row and column representing a
vertex, and elements where corresponding vertices for the row-column pair are
adjacent equals 1, the rest of the elements equal 0. Let $x$ be the vector of
the values $x_v$. Each element of $Ax$ is the number of colored vertices in
the neighborhood of the corresponding vertex, so if $T$ is an efficient open
dominating set then
\[
  Ax = \mathbf{1},
\]
where $\mathbf{1}$ is the vector with every element equal to 1.

\section{The odd graph $K(2k+1,\,k)$}

% CHANGED: this is the claim from the first sentence of page 1 of the notes,
% with "efficient dominating set" written "efficient open dominating set".
The odd graph $K(2k+1,\,k)$ has no efficient open dominating set.

All Kneser graphs are regular, so every vertex has the same degree. For our
$K(2k+1,\,k)$ graph:
\[
  \binom{2k+1-k}{k} \;=\; \binom{k+1}{k} \;=\; \frac{(k+1)!}{1!\,k!} \;=\; k+1
\]

If we have sets of size $k$, then we have $(2k+1)-k$ sets to form adjacent
vertices with.

Thus, each row of the adjacency matrix has $k+1$ 1's. We can try to solve the
matrix equation, ignoring coloring rules for now.

$\rightarrow$ We want every value in the product vector to be a 1, so every
value of $x$ has to be the same.

% CHANGED: the right-hand side "1" is written \mathbf{1} (all-ones vector),
% here and in every matrix equation below.
\[
  A \cdot x \;=\; A \cdot
  \begin{bmatrix} c \\ c \\ \vdots \\ c \end{bmatrix}
  \;=\; \mathbf{1}
\]

When we multiply each row of $A$, we get:
\[
  \underbrace{c + c + c + \cdots + c}_{k+1} \;=\; (k+1) \cdot c
\]

Each row has to equal 1, so:
\[
  (k+1) \cdot c = 1, \qquad c = \frac{1}{k+1}
\]

So
\[
  x \;=\;
  \begin{bmatrix} 1/(k+1) \\ 1/(k+1) \\ \vdots \\ 1/(k+1) \end{bmatrix}
  \;=\; \frac{1}{k+1}
  \begin{bmatrix} 1 \\ 1 \\ \vdots \\ 1 \end{bmatrix}
\]

Now we can find the eigenvalues of our $Ax = \mathbf{1}$ matrix using the odd
graph formula
\[
  \lambda_j = (-1)^j\,(k+1-j), \qquad j = 0, 1, \dots, k
\]

% CHANGED: "the matrix in invertible" -> "is invertible".
The range of magnitudes for our eigenvalue is from $1$ ($k = j$) to $k+1$
($0 = j$). $0$ is not in this range, so the matrix is invertible.

% CHANGED: the exponent n in N_0^n is written |V(K)|, the number of vertices,
% so that it does not collide with the n of K(2k+n, k).
Now suppose there was some vector $w$ such that $Aw = \mathbf{1}$, and
$w \in \mathbb{N}_0^{|V(K)|}$ such that it meets CFON$^{*}$ coloring
restrictions.

Then:
\begin{align*}
  Aw &= Ax = \mathbf{1} \\
  A^{-1}Aw &= A^{-1}Ax && \text{(we've proven invertibility)} \\
  w &= x,
\end{align*}
% CHANGED: "x is fraction" -> "x is fractional".
but $x$ is fractional which contradicts the definition of $w$. Thus, $w$
cannot exist, and $x$ is the unique solution.

% CHANGED: "it's" -> "its" (twice); "not efficient dominating set exists" ->
% "no efficient open dominating set exists".
Since $x$ is the unique solution and its elements are fractional, and since
$x_v$ can only equal 0 or 1 based on its inclusion in $T$, no efficient open
dominating set exists and thus $K(2k+1,\,k)$ is not 1-colorable.

Thus, $K(29,\,14)$ is not 1-colorable.

\section{$K(2k+2,\,k)$}

Now for the existence of an efficient open dominating set for $K(2k+2,\,k)$.

The $K(2k+2,\,k)$ graph is $\dbinom{k+2}{k}$-regular.

% CHANGED: the notes have k^2+3k+2. The binomial coefficient is
% (k+2)(k+1)/2 = (k^2+3k+2)/2, so the factor of 2 is restored here and in the
% three places below that use it.
So, we have $\dfrac{k^2+3k+2}{2}$ degrees for every vertex.

Let $B$ be the adjacency matrix, with each row and column representing a
vertex, and elements where corresponding vertices for the row-column pair are
adjacent equals 1, the rest of the elements equal 0.

Note again that this matches with the values assigned based on $T$:
\[
  x_v =
  \begin{cases}
    1, & v \in T\\
    0, & v \notin T
  \end{cases}
\]

% CHANGED: factor of 2 (see above).
Each row of the adjacency matrix has $\dfrac{k^2+3k+2}{2}$ 1's.

Then, suppose
\[
  By \;=\; \begin{bmatrix} 1 \\ 1 \\ \vdots \\ 1 \end{bmatrix}
\]

Suppose all the values of $y$ are some constant $d$. Then:
% CHANGED: factor of 2 (see above). The notes have (k^2+3k+2) d = 1 and
% d = 1/(k^2+3k+2).
\[
  \left( \frac{k^2+3k+2}{2} \right) d = 1
\]
\[
  d = \frac{2}{k^2+3k+2} < 1, \qquad \forall k,\; k \in \mathbb{Z}^{+}
\]

The eigenvalues for $B$:
\[
  \lambda_j = (-1)^j \binom{k+2-j}{2} \qquad \text{for } j = 0, 1, \dots, k
\]

If $j = 0$, then $\dbinom{k+2}{2} > \dbinom{2}{2} = 1 > 0$.

If $j = k$, then $\dbinom{2}{2} = 1 > 0$.

Thus the eigenvalue for $B$ is always non-zero, making $B$ invertible.

% CHANGED: "BWOC" is spelled out; N_0^n -> N_0^{|V(K)|} as in Section 1.
By way of contradiction, assume a matrix $p$ exists where $Bp = \mathbf{1}$,
and $p \in \mathbb{N}_0^{|V(K)|}$.

\[
  Bp = By, \qquad B^{-1}Bp = B^{-1}By, \qquad p = y,
\]
which is a contradiction because $y$ is fractional and $p$ can't be. Thus
there is no efficient open dominating set for $K(2k+2,\,k)$.

\end{document}
